% ============================================================
% Auto-generated by harryopo-convert
% Compile: cd templates && .\build.ps1
% ============================================================
\documentclass[nomath]{harryopo-report}

\title{智慧交通流量预测研究}
\author{张三、李四}
\date{2026年06月22日}

\begin{document}

\maketitle

\begin{abstract}
本文提出了一种基于深度学习的交通流量预测方法，通过融合LSTM与注意力机制，在真实数据集上取得了优于传统方法的预测精度。
\end{abstract}

\tableofcontents

\introduction
本文由 harryopo-convert 自动生成。

\section{引言}

随着城市化进程加快，交通拥堵问题日益严重。准确的交通流量预测对于智能交通系统至关重要。近年来，深度学习在时序预测领域取得了显著成果。

本文的主要贡献如下：

\begin{itemize}
  \item 提出了一种融合LSTM与注意力机制的交通预测模型
  \item 在三个真实数据集上进行了全面验证
  \item 模型计算效率提升约30%
\end{itemize}

\subsection{相关背景}

传统方法主要依赖ARIMA等统计模型，但这些方法难以捕捉复杂的非线性关系。深度学习方法通过多层非线性变换，能够更有效地学习交通流的时空特征。

\section{方法}

\subsection{数据预处理}

我们采用了以下三阶段流程：

\begin{enumerate}
  \item 缺失值填充：使用线性插值填补传感器数据缺失
  \item 归一化处理：采用Min-Max归一化至[0,1]区间
  \item 时间窗口切片：构造长度为24的滑动窗口
\end{enumerate}

\subsection{模型架构}

核心模型基于LSTM编码器-解码器架构：

\begin{lstlisting}[style=pystyle,caption={}]
class TrafficModel(nn.Module):
    def __init__(self, input_dim=64, hidden_dim=128, num_layers=2):
        super().__init__()
        self.lstm = nn.LSTM(input_dim, hidden_dim, num_layers)
        self.attention = nn.MultiheadAttention(hidden_dim, 8)
        self.fc = nn.Linear(hidden_dim, 1)

    def forward(self, x):
        lstm_out, _ = self.lstm(x)
        attn_out, _ = self.attention(lstm_out, lstm_out, lstm_out)
        return self.fc(attn_out[-1])
\end{lstlisting}

模型训练使用Adam优化器，初始学习率为0.001。损失函数采用均方误差（MSE）。

\section{实验}

\subsection{数据集与评估指标}

我们在三个公开数据集上进行实验：

\begin{table}[htbp]
  \centering
  \caption{}
  \begin{tabular}{l l l l}
    \toprule
    数据集 & 节点数 & 时间跨度 & 采样频率 \\
    \midrule
    PeMSD4 & 307 & 59天 & 5分钟 \\
    PeMSD8 & 170 & 62天 & 5分钟 \\
    METR-LA & 207 & 121天 & 5分钟 \\
    \bottomrule
  \end{tabular}
\end{table}

\subsection{主要结果}

\begin{quote}
实验结果表明，我们的模型在所有数据集上均优于基线方法，尤其在长期预测（60分钟）场景下优势更加显著。
\end{quote}

模型的损失函数定义为：

\begin{equation}
\mathcal{L} = \frac{1}{N} \sum_{i=1}^{N} (y_i - \hat{y}_i)^2 + \lambda \|\theta\|_2
\end{equation}

其中 $\lambda=0.001$ 为L2正则化系数。

\section{结论}

本文提出了一个基于LSTM-注意力机制的交通流量预测模型。实验结果表明，该模型在三个公开数据集上均取得了最优性能。

\begin{thebibliography}{99}
\bibitem{ref1} [1] Wu et al. Graph WaveNet for Deep Spatial-Temporal Graph Modeling. IJCAI 2019.
\bibitem{ref2} [2] Li et al. Diffusion Convolutional Recurrent Neural Network. ICLR 2018.
\bibitem{ref3} [3] Bai et al. Adaptive Graph Convolutional Recurrent Network. NeurIPS 2020.
\end{thebibliography}

\end{document}